% =====================================================================
% Literature Review (revised)
% Restructured to follow the thesis pipeline:
%   seasonal wind-gust motivation -> marginal tail modelling and
%   standardisation -> extremal dependence and extremal copulas
%   (pairwise summaries) -> temporal clustering/declustering ->
%   max-stable processes -> Brown--Resnick -> pairwise composite
%   likelihood -> block resampling and boundary problems in testing ->
%   simulation benchmark -> positioning.
% Each subsection answers: what problem does this literature solve;
% why does this thesis need it; what object/theorem is introduced;
% what does it allow us to estimate or diagnose; what are its
% limitations here.
% The F-madogram is deliberately demoted to a single remark.
% New citation keys (BibTeX supplied separately): sklar1959,
% gudendorf2010, schlather2003, ledford1996, leadbetter1983, ferro2003,
% dehaan1984, husler1989, lindsay1988, efron1979, davisonhinkley1997,
% chernoff1954, self1987, andrews2000, andrews2001.
% Keys reused from the previous draft are assumed to be in the .bib:
% fisher1928, gnedenko1943, pickands1975, balkema1974, dehaan2006,
% coles2001, coles1999, davison2012, davison2013, smith1990,
% schlather2002, brown1977, kabluchko2009, padoan2010, huser2013,
% huserdavison2014, varin2011, cooley2006, dombry2016, buhl2016,
% kuensch1989, lahiri2003, oesting2017, sharkey2020, pacey2021,
% gliksman2023european, huseropitzthibaud2017, huserwadsworth2019,
% huserwadsworth2022.
% =====================================================================
\section{Literature Review}\label{sec:literature}

This chapter assembles the statistical literature on which the analysis rests, in the order in which the analysis uses it. The empirical question --- whether winter wind-gust extremes in the Netherlands have a wider spatial footprint than summer ones --- is answered through a chain of steps: marginal tails are modelled and standardised to a common scale; temporal clustering of stormy days is diagnosed; extremal dependence between station pairs is summarised nonparametrically through extremal copulas; a Brown--Resnick max-stable model pools those pairwise summaries into a two-parameter curve per season; the parameters are estimated by pairwise composite likelihood; uncertainty is quantified by a season-year block bootstrap; and the whole procedure is benchmarked by Monte Carlo simulation under a known data-generating process. Each subsection below follows the same logic: what problem the literature solves, why this thesis needs it, which object or theorem it introduces, what it allows the thesis to estimate or diagnose, and where its limitations lie. The aim is that the reader can verify, link by link, why each component of the methodology in Section~\ref{sec:methodology} is present.

\subsection{Wind-gust extremes and the seasonal footprint question}\label{sec:lit-wind}

Extreme wind gusts are among the costliest weather hazards in north-western Europe, and much of the cost arises because a severe storm strikes many locations at once, so that losses aggregate over space. For European extratropical windstorms, \citet{sharkey2020} model the spatial extent and the severity of an event jointly and argue that the footprint of a storm --- not only its peak intensity at a single site --- is a primary object of inference for aggregate risk. Their work establishes the perspective adopted here: the spatial extent of an extreme event is itself an estimand, distinct from the marginal severity at any one station.

The seasonal contrast is physically motivated. Winter wind extremes in the Netherlands are typically produced by synoptic-scale extratropical cyclones, whose mature wind fields span hundreds of kilometres, whereas summer extremes are more often produced by convective storms and downbursts, which are intense but spatially confined. \citet{pacey2021} document in a European climatology of severe convective windstorms that such events are concentrated in the warm season, and \citet{gliksman2023european} review how both storm types translate into wind damage. The meteorological contrast therefore suggests a statistical hypothesis that is not about marginal magnitudes --- a winter gust and a summer gust at one station may look alike --- but about dependence: winter extremes should co-occur over larger distances than summer extremes. The estimand of this thesis is a seasonal distance scale of extremal dependence, and the remainder of this chapter builds the tools needed to define, estimate, and test it.

\subsection{Extreme-value theory for marginal tail modelling}\label{sec:lit-evt}

The first problem is marginal: statements about extreme gusts require a model for the upper tail of each station's gust distribution, in a region where empirical frequencies are sparse and bulk-fitted models are unreliable. Extreme-value theory supplies the limiting families for this purpose. For maxima \(M_n\) of independent and identically distributed variables, the Fisher--Tippett--Gnedenko theorem \citep{fisher1928,gnedenko1943} states that if \((M_n-b_n)/a_n\) converges to a non-degenerate law for some normalising sequences, that law must be a generalised extreme value (GEV) distribution, with location \(\mu\), scale \(\sigma\), and a shape \(\xi\) that governs whether the tail is bounded (\(\xi<0\)), light (\(\xi=0\)), or heavy (\(\xi>0\)); \citet{dehaan2006} give the modern treatment and \citet{coles2001} the standard applied account. A companion route models threshold exceedances: by the Pickands--Balkema--de Haan theorem \citep{pickands1975,balkema1974}, exceedances over a high threshold are approximately generalised Pareto precisely when the maxima are in a GEV domain of attraction, so block-maximum and peaks-over-threshold analyses are two views of the same tail.

This thesis needs marginal extreme-value modelling for two reasons. First, the spatial dependence model of Section~\ref{sec:lit-br} is specified on a fixed common marginal scale, conventionally unit Fr\'echet \citep{dehaan2006}, so each station--season margin must be transformed onto that scale before dependence can be estimated. Second, the margins genuinely differ: coastal stations are windier than inland ones, and winter is windier than summer at the same station. Without standardisation, a winter--summer ``dependence'' comparison would partly reflect these marginal climate differences rather than co-occurrence of extremes. The operational device is the probability integral transform: a parametric distribution \(G_i^{(c)}\) is fitted to the daily maxima of station \(i\) in season \(c\), and \(U=G_i^{(c)}(X)\) (uniform margins) and \(Z=-1/\log U\) (unit-Fr\'echet margins) place all station--season samples on the same scale, exactly as implemented in Section~\ref{sec:meth-marginals}.

The role of the GEV here must be stated with care. A daily maximum gust is the maximum over many sub-daily measurements within one day, which makes an extreme-value family a natural candidate for its distribution, but a one-day block is short and the asymptotic block-maximum regime is not literally verified. The GEV family is therefore used as a parametric upper-tail standardisation model: a flexible three-parameter family whose adequacy is a matter of fit, not a claim that daily maxima are exact asymptotic block maxima. This distinction matters because the standardisation is the hinge between marginal modelling and dependence modelling: any marginal misspecification propagates into the transformed fields \(U\) and \(Z\) and hence into every dependence summary computed from them. Two safeguards follow. Marginal fit is checked station by station, and a threshold-based alternative exists --- generalised Pareto tails combined with a censored pairwise likelihood that fits dependence only above high levels \citep{huser2013} --- which is noted again in Section~\ref{sec:lit-cl} as the natural robustness extension. After the transform, remaining winter--summer differences are, by construction, differences in dependence; that separation is made precise by Sklar's theorem in the next subsection.

\subsection{From marginal extremes to extremal dependence}\label{sec:lit-dependence}\label{sec:lit-diagnostics}

The second problem is dependence: how do extremes at two stations co-occur, and how does that co-occurrence decay with distance? Ordinary correlation does not answer this. It is a single second-moment summary dominated by the bulk of the data --- the many ordinary windy days --- and two bivariate distributions can share the same correlation while one places substantial mass on joint extremes and the other places almost none \citep{coles1999}. A dependence concept specific to the joint tail is required.

Copulas provide the language for separating dependence from margins. By Sklar's theorem \citep{sklar1959}, a continuous joint distribution factorises into its margins and a copula \(C\), so that once margins are standardised as in Section~\ref{sec:lit-evt}, all questions about joint extremes are questions about \(C\) near its upper corner. For componentwise maxima the limiting copulas form a special class, the extreme-value (or extremal) copulas \citep{dehaan2006,gudendorf2010}. In the bivariate case they can be written through the exponent function \(V\) used throughout this thesis,
\[
    C(u_1,u_2)=\exp\!\Bigl\{-V\Bigl(-\tfrac{1}{\log u_1},\,-\tfrac{1}{\log u_2}\Bigr)\Bigr\},
\]
with \(V\) homogeneous of order \(-1\); equivalently through the Pickands dependence function \(A\) on \([0,1]\). Extreme-value copulas are exactly the dependence structures compatible with max-stability, so they are also the finite-dimensional dependence structures of the spatial models introduced in Section~\ref{sec:lit-maxstable}: the nonparametric and parametric strands of this thesis describe the same object.

Two scalar functionals of \(C\) carry the pairwise analysis. The first is the extremal coefficient \(\theta=V(1,1)\in[1,2]\), which satisfies \(C(u,u)=u^{\theta}\) on the diagonal and is read as the effective number of independent stations in the pair: \(\theta=1\) is complete dependence and \(\theta=2\) independence; its properties and estimation are studied by \citet{schlather2003}, and the derivation is collected in Appendix~\ref{app:exponent}. The second is the finite-level tail-dependence coefficient \(\chi_u=\Pr(U_j>u\mid U_i>u)\), the conditional probability that one station is extreme given that the other is; under an extreme-value copula \(\chi_u=(1-2u+u^{\theta})/(1-u)\), with limit \(\chi=\lim_{u\uparrow1}\chi_u=2-\theta\) \citep{coles1999}, as derived in Appendix~\ref{app:tail-dependence}. Both quantities are transparent in the sense the introduction requires: they can be computed for every station pair from the uniform-scale data alone, without any spatial model, and plotted against the inter-station distance \(d_{ij}\). This thesis uses them in two roles --- first as model-free evidence that seasonal dependence differs, and later as diagnostics for the parametric fit (Section~\ref{sec:meth-diagnostics}). Nonparametric estimates of \(\theta\) are available directly from the empirical copula evaluated on the diagonal at high levels; the same coefficient can also be reached through the \(F\)-madogram of \citet{cooley2006}, a first-order variogram of the uniform-scale field, but since that route estimates the same target with additional geostatistical machinery, it is mentioned here only for completeness and is not developed further in this thesis.

One caveat from this literature shapes the diagnostics. The limit \(\chi=0\) defines asymptotic independence, and a pair can be asymptotically independent while exhibiting substantial dependence at every finite level; \citet{ledford1996} introduce a residual tail-dependence framework that quantifies how \(\chi_u\) decays to zero in that case. The caveat matters because the max-stable models of Section~\ref{sec:lit-maxstable} impose \(\chi(d)>0\) wherever \(\theta(d)<2\), so if gusts are in fact asymptotically independent, a fitted max-stable curve will overstate joint exceedances at the rarest levels \citep{davison2013,huserwadsworth2022}. The empirical \(\widehat\chi_u\) therefore has to be inspected for decay across levels and distances faster than the fitted model allows. Finally, the pairwise viewpoint is adopted not only because it is interpretable: with \(N=33\) stations and a few thousand days per season, the joint tail in \(33\) dimensions cannot be estimated nonparametrically with any useful precision, whereas pairwise objects indexed by distance can be. The parametric model exists precisely to pool the \(528\) noisy pairwise summaries into a smooth two-parameter curve per season, and the same pairwise structure reappears in the likelihood of Section~\ref{sec:lit-cl}.

\subsection{Temporal clustering and declustering of extremes}\label{sec:lit-temporal}

Everything above treats observations as independent replicates, but daily gust maxima at one station form a time series, and storms persist: a single windstorm can produce extreme daily maxima on consecutive days, so extreme days arrive in clusters. The theory of extremes for stationary sequences describes the consequence. Under a weak long-range mixing condition, the maximum of \(n\) dependent observations still has a GEV limit, but modified by the extremal index \(\theta_{\mathrm T}\in(0,1]\), in the sense that \(\Pr(M_n\le u_n)\approx F(u_n)^{n\theta_{\mathrm T}}\) \citep{leadbetter1983,coles2001}; the notation \(\theta_{\mathrm T}\) is used here to avoid confusion with the spatial extremal coefficient \(\theta\). Values \(\theta_{\mathrm T}<1\) indicate clustering, \(1/\theta_{\mathrm T}\) approximates the mean cluster size, and a series of \(n\) days behaves, for its extremes, like roughly \(n\theta_{\mathrm T}\) independent days. Estimation and cluster identification are well developed: runs declustering merges exceedances separated by fewer than a chosen run length into one cluster \citep{coles2001}, at the cost of sensitivity to that tuning choice, and the intervals estimator of \citet{ferro2003} estimates \(\theta_{\mathrm T}\) from inter-exceedance times and yields an automatic declustering without a run parameter.

This thesis needs the clustering literature for three reasons, none of which is to decluster the data before estimation. First, the marginal GEV standardisation of Section~\ref{sec:lit-evt} remains a fit of the distribution of daily maxima --- clustering does not bias that fitted distribution under stationarity --- but it would invalidate independence-based uncertainty statements about the marginal parameters. Second, the pairwise composite likelihood of Section~\ref{sec:lit-cl} treats days as replicates of the spatial field; serial clustering leaves its estimating equations valid under stationarity and mixing \citep{huserdavison2014}, but inflates the variability of the composite score, so uncertainty must not be computed as if days were independent. Third, any resampling scheme must keep the days of one storm together, which dictates the block structure of Section~\ref{sec:lit-inference}. Accordingly, temporal clustering enters the thesis as a pre-inference diagnostic --- the extremal index is estimated per station and season to gauge how strongly extreme days cluster and whether winter and summer differ in persistence --- while the inferential consequences are absorbed by resampling whole season-years. Declustering the data themselves before the spatial fit is deliberately avoided: removing within-cluster days would discard entire spatial fields that carry information about cross-station dependence and would create an irregular panel. The limitation of this choice is the mirror image of its advantage: a diagnostic quantifies serial dependence but does not remove it, so the validity of the uncertainty statements rests on the block resampling that the diagnostic motivates.

\subsection{Spatial extremes and max-stable processes}\label{sec:lit-maxstable}

The pairwise summaries of Section~\ref{sec:lit-dependence} describe each station pair separately; the next problem is to model all pairs coherently as one spatial object. Max-stable processes are the spatial analogue of the GEV limit: if pointwise maxima of independent copies of a random field, suitably normalised, converge to a non-degenerate limit field, that limit must be max-stable \citep{dehaan1984,dehaan2006}. The spectral representation of \citet{dehaan1984} expresses such a process as the pointwise maximum over a Poisson ensemble of random storm profiles, which is both the source of concrete models and the basis of the exact simulation algorithms used later \citep{dombry2016}. Crucially for this thesis, the finite-dimensional distributions of a max-stable process have GEV margins joined by extreme-value copulas, so each pair of sites carries an extremal coefficient \(\theta(d)\) and a tail-dependence curve \(\chi_u(d)\); under stationarity and isotropy these depend on the pair only through its distance. A stationary isotropic max-stable model is thus the pooling device the analysis needs: it asserts that all \(528\) station pairs share one dependence-versus-distance curve per season, against which the pairwise clouds of Section~\ref{sec:lit-dependence} can be checked.

Several parametric families exist, and their known deficiencies explain the model choice. The model of \citet{smith1990} builds storms with deterministic profiles and implies extremal surfaces that are unrealistically smooth and rigid. The model of \citet{schlather2002} draws storm profiles from a stationary Gaussian process and is more flexible, but its extremal coefficient is bounded away from \(2\) --- with a vanishing underlying correlation it cannot exceed \(1+\sqrt{1/2}\approx1.71\) --- so it cannot represent independence between distant stations, which a country-spanning gust network plausibly approaches. \citet{davison2012} review these families, recommend the Brown--Resnick process as the leading stationary model, and document a point that motivates the entire design of this thesis: a model can fit every margin well and still misrepresent joint extremes, which is why dependence is modelled directly here rather than inferred from marginal fits. The limitation at the level of the whole class is the one already raised: max-stability is an asymptotic idealisation, adopted as an approximation for daily fields, and its finite-level adequacy is an empirical question for the diagnostics, not an assumption.

\subsection{The Brown--Resnick process}\label{sec:lit-br}

The Brown--Resnick process originates with \citet{brown1977}, who constructed a max-stable process from Brownian motion with drift, and was generalised by \citet{kabluchko2009}, who build max-stable fields from log-Gaussian spectral functions based on an intrinsically stationary Gaussian process \(W\) and characterise when the result is stationary. Two features of their characterisation make the family ideal for this thesis. First, the law of the max-stable field depends on \(W\) only through its variogram \(\gamma\), so the variogram is the complete dependence parameter. Second, the bivariate margins are of H\"usler--Reiss type \citep{husler1989}, with the closed-form exponent \(V_d\) stated in Section~\ref{sec:meth-variogram}, so the implied extremal coefficient at distance \(d\) is the explicit curve \(\theta(d)=2\Phi\bigl(\sqrt{\gamma(d)}/2\bigr)\).

With the isotropic powered-exponential variogram \(\gamma(d)=(d/\rho)^{\alpha}\), \(\rho>0\), \(\alpha\in(0,2]\), the dependence structure is reduced to two interpretable parameters. The range \(\rho\) is a distance scale in kilometres: at \(d=\rho\) the variogram equals one, so \(\theta(\rho)=2\Phi(1/2)\approx1.38\) for every \(\alpha\), making \(\rho\) the distance at which a fixed reference level of extremal dependence is reached. The smoothness \(\alpha\) controls how steeply \(\theta(d)\) rises from \(1\) at \(d=0\) towards \(2\) at long range. The seasonal research question thereby becomes a parametric comparison --- is \(\rho_{\mathrm W}>\rho_{\mathrm S}\)? --- read jointly with \(\alpha_c\) and the implied curves \(\theta_c(d)\), exactly as formalised in Section~\ref{sec:meth-comparison}. (Variogram conventions differ across references by a factor of two; the convention used throughout this thesis is fixed in Section~\ref{sec:meth-variogram}.) The family is also a proven choice for this data type: beyond its status as the standard stationary model for spatial extremes \citep{davison2012}, \citet{oesting2017} apply bivariate Brown--Resnick processes to extreme wind gusts over a station network in northern Germany, evidence that the dependence structure is suitable for north-west European gust fields, and anisotropic space--time extensions exist \citep{buhl2016} should the simplifications below need relaxing.

Three limitations qualify the model in this application. First, isotropy: dependence is assumed to depend on distance alone, whereas the Dutch coastline and prevailing storm tracks could induce directional effects; isotropy is retained for parsimony and identifiability on a \(33\)-station network and is flagged as an extension rather than tested. Second, asymptotic dependence: Brown--Resnick fields satisfy \(\chi(d)=2-\theta(d)>0\) at every finite distance, so if gust extremes are asymptotically independent the model overstates the co-occurrence of the very rarest events \citep{davison2013}; flexible alternatives that bridge or estimate the dependence class exist, including Gaussian scale mixtures \citep{huseropitzthibaud2017}, models with unknown extremal dependence class \citep{huserwadsworth2019}, and the broader toolbox reviewed by \citet{huserwadsworth2022}. Third, the process is a model, not the data-generating mechanism: it is used as a parsimonious approximation whose adequacy at the finite levels and distances the network resolves is examined by the diagnostics of Section~\ref{sec:meth-diagnostics}, and the seasonal conclusions are confined to those levels and distances rather than extrapolated to rarer events.

\subsection{Pairwise composite likelihood}\label{sec:lit-cl}

Having a model does not yet give an estimator, because full-likelihood inference for max-stable processes is infeasible beyond a handful of sites. The joint distribution function at \(N\) sites is \(\exp\{-V(z_1,\dots,z_N)\}\), and differentiating to obtain the density produces a sum over all partitions of the \(N\) sites; the number of terms is the Bell number, which already exceeds \(10^5\) for \(N=10\) and is astronomically large for the \(N=33\) stations used here \citep{padoan2010}. A principled surrogate objective is required.

Composite likelihood supplies it. The idea, formalised by \citet{lindsay1988} and reviewed by \citet{varin2011}, is to multiply valid low-dimensional likelihood components in place of the intractable joint density; the pairwise version used in spatial extremes is
\[
    \ell(\vartheta)=\sum_{i<j}\sum_{t}\log f_{ij}\bigl(z_{i,t},z_{j,t};\vartheta\bigr),
\]
where \(f_{ij}\) is the exact bivariate density implied by the model. Because each \(f_{ij}\) is a true density, the composite score has mean zero at the true parameter, so the estimating equations are unbiased and, under regularity conditions and independent replicates, the maximiser is consistent and asymptotically normal with covariance given by the inverse Godambe information,
\[
    G(\vartheta)^{-1}=H(\vartheta)^{-1}J(\vartheta)\,H(\vartheta)^{-1},
\]
with \(H\) the sensitivity matrix (the expected negative Hessian of the composite log-likelihood) and \(J\) the variability matrix (the variance of the composite score) \citep{varin2011,padoan2010}. The precise sense in which ``ordinary standard errors do not apply'' is that the second Bartlett identity \(H=J\) holds for a genuine likelihood but fails for a composite one, so the inverse Hessian is not the asymptotic variance and a sandwich correction --- or a resampling substitute --- is mandatory. The price of the surrogate is a loss of efficiency relative to the infeasible full-likelihood estimator; the gains are tractability and the fact that only the bivariate margins, which the Brown--Resnick model specifies in closed form, need to be correct.

For max-stable processes specifically, \citet{padoan2010} develop the pairwise likelihood and its asymptotic theory with independent temporal replicates, and \citet{huser2013} treat the Brown--Resnick case, including a censored pairwise likelihood that uses only pairs above a high threshold so that the fit targets the joint tail --- the natural companion to the threshold-based marginal route noted in Section~\ref{sec:lit-evt}. \citet{huserdavison2014} extend pairwise likelihood to space--time extremes and show that temporally dependent replicates can be accommodated under mixing conditions, which is the theoretical licence for applying the estimator to daily fields that cluster in time. Pair selection and weighting can trade efficiency against computation; with only \(528\) pairs all of them are retained here, and the long-distance pairs are in fact informative about the range \(\rho\), as discussed in Section~\ref{sec:meth-cl}. The practical consequence of this subsection for inference is the status of \(J\): with independent daily replicates it is estimable by the empirical variance of daily score contributions, but with serially clustered days (Section~\ref{sec:lit-temporal}) those contributions are dependent and the naive estimate is biased downward, typically badly so. Valid options are window-type corrections or resampling that respects the dependence; this thesis takes the resampling route developed next, retaining the independence-based Godambe form only as a theoretical reference, as stated in Section~\ref{sec:meth-inference}.

\subsection{Resampling for dependent data and boundary problems in testing}\label{sec:lit-inference}

The bootstrap of \citet{efron1979} estimates sampling uncertainty by resampling individual observations with replacement, and is valid for independent data and smooth statistics. Applied naively to the present data it fails twice. Resampling individual days treats serially dependent storm days as exchangeable, destroying the time dimension of the data; for positively dependent series this typically understates the variability of estimators and produces intervals that are too narrow. And any scheme that resampled station--days separately would in addition destroy the same-day cross-station dependence --- which is not a nuisance here but the estimand itself. A valid resampling scheme must therefore keep whole daily fields intact and keep temporally adjacent fields together.

The block bootstrap is the established solution for dependent data. \citet{kuensch1989} introduced the moving-block bootstrap for stationary sequences: resampling contiguous blocks preserves the dependence within each block, and for smooth functionals of stationary mixing series the procedure is consistent when the block length grows suitably with the sample size, with a bias--variance trade-off in the choice of block length; \citet{lahiri2003} gives a comprehensive treatment of these methods and their theory, and \citet{davisonhinkley1997} the standard practice for percentile intervals. The present data offer a natural block that removes the block-length tuning problem: the season-year, one whole winter or one whole summer. Because each winter is assembled as a contiguous December--February block (Section~\ref{sec:data}), no storm is split across blocks; because entire days are resampled as \(33\)-station fields, the spatial dependence within a day is untouched; and because winters (or summers) of different years are separated by months, treating blocks as approximately independent is plausible at the storm timescale --- with the acknowledged caveat that low-frequency climate variability can correlate blocks across years, as discussed in Section~\ref{sec:meth-comparison}. Re-estimating the marginal transform within every bootstrap replication, as Section~\ref{sec:meth-inference} prescribes, additionally propagates the standardisation uncertainty of Section~\ref{sec:lit-evt} into the intervals. The cost of the natural block is its scarcity: with only about \(35\) season-years per season, the bootstrap distribution is built from few blocks, percentile intervals are correspondingly coarse, and this finite-block resolution is one of the properties the simulation study is designed to quantify.

\paragraph{Boundary problems in testing dependence.}
A final inferential point from the literature shapes the form of the hypothesis test. A seemingly natural null --- ``no extremal dependence'', expressed as \(\chi=0\), \(\theta(d)=2\), or a dependence parameter at the edge of its space --- is statistically nonstandard, because the null value then lies on the boundary of the parameter space: likelihood-ratio and Wald statistics no longer have their usual chi-squared limits but mixtures thereof \citep{chernoff1954,self1987}, the general asymptotic theory is substantially more involved \citep{andrews2001}, and even the bootstrap can be inconsistent when the true parameter sits on a boundary \citep{andrews2000}, so resampling does not repair the problem. In the parametrisation used here the issue is concrete: independence at all observed distances corresponds to the lower edge of the range space, since \(\theta_c(d)\to2\) as \(\rho_c\downarrow0\), so a test of zero spatial dependence would place its null exactly where the standard theory breaks down. This thesis therefore follows the comparison route instead: the hypothesis of Section~\ref{sec:meth-comparison} is the equality of the two seasonal ranges, \(H_0:\rho_{\mathrm W}=\rho_{\mathrm S}\) against \(H_1:\rho_{\mathrm W}>\rho_{\mathrm S}\), a contrast between two interior parameters estimated from disjoint seasonal samples, for which standard block-bootstrap reasoning applies. Boundary effects can still arise mechanically --- the smoothness is bounded by \(\alpha\le2\) and the optimiser works on a box --- and, in the spirit of \citet{andrews2001}, these are monitored rather than tested: the simulation records boundary-hit rates, and estimates piling up at an edge are read as weak identification rather than fed into a formal test.

\subsection{Simulation as a controlled benchmark}\label{sec:lit-simulation}

In the empirical data the true dependence structure is unknown, so estimator bias, optimiser artefacts, and miscalibration of the test cannot be distinguished from genuine seasonal signal by looking at the data alone. The standard remedy is a Monte Carlo benchmark: simulate from a Brown--Resnick process with known \((\rho_c,\alpha_c)\) on the observed \(33\)-station geometry and with the empirical season-year layout, apply the entire pipeline unchanged, and ask whether it recovers what was put in. This validates the procedure under controlled conditions --- bias and root-mean-squared error of \(\widehat\rho_c\), \(\widehat\alpha_c\) and \(\widehat\Delta_\rho\); coverage of the bootstrap intervals; size of the one-sided test when \(\Delta_\rho=0\) and power when \(\Delta_\rho>0\) --- before the same procedure is trusted on the gust data, and it connects directly to the seasonal comparison because the alternative scenario plants exactly the kind of range difference the empirical analysis seeks. Exactness of the simulator matters: approximate max-stable simulators can distort the joint tail that is under study, and the algorithms of \citet{dombry2016} simulate Brown--Resnick processes exactly via their extremal functions, so the data-generating process in the benchmark is the model itself and not an approximation to it. The limits of the exercise are stated as clearly as its purpose: the benchmark is conditional on the assumed model, so it cannot certify that observed gusts are Brown--Resnick --- that question belongs to the diagnostics --- and because the simulated daily fields are independent, it exercises the block bootstrap without genuine within-block dependence, making the reported calibration a best case, as discussed in Section~\ref{sec:simulation}.

\subsection{Positioning of this thesis}\label{sec:lit-seasonal}

Allowing dependence to differ between winter and summer is the simplest form of non-stationary extremal dependence: season acts as a binary covariate, and one Brown--Resnick model is estimated per season. The alternative of fitting a single model to all days would mix synoptic winter storms with convective summer events and blur exactly the contrast of interest, while richer designs --- covariate-dependent dependence structures or anisotropic space--time max-stable models \citep{buhl2016,huserdavison2014} --- are natural extensions that this thesis deliberately does not pursue, both because a two-season contrast does not require them and because their additional parameters would be harder to identify from \(33\) stations.

This thesis develops no new estimator, model, or theorem. Its contribution is to combine well-understood components into one transparent and validated workflow for a specific empirical question: GEV-based upper-tail standardisation of the margins \citep{fisher1928,gnedenko1943,coles2001}; an explicit temporal-clustering diagnostic before spatial inference \citep{ferro2003}; nonparametric pairwise extremal-dependence evidence rooted in extreme-value copulas \citep{coles1999,schlather2003}; a Brown--Resnick model estimated by pairwise composite likelihood \citep{kabluchko2009,padoan2010,huser2013}; a season-year block bootstrap for dependent data \citep{kuensch1989,lahiri2003}; a hypothesis framed as an interior seasonal comparison rather than a boundary test \citep{andrews2001}; and an exact-simulation benchmark of the whole procedure \citep{dombry2016}. Each component answers a failure mode of a simpler approach, and together they allow the seasonal footprint question of Section~\ref{sec:lit-wind} to be answered with stated, checkable assumptions.